%%=============================================================================
%% Samenvatting
%%=============================================================================

% TODO: De "abstract" of samenvatting is een kernachtige (~ 1 blz. voor een
% thesis) synthese van het document.
%
% Een goede abstract biedt een kernachtig antwoord op volgende vragen:
%
% 1. Waarover gaat de bachelorproef?
% 2. Waarom heb je er over geschreven?
% 3. Hoe heb je het onderzoek uitgevoerd?
% 4. Wat waren de resultaten? Wat blijkt uit je onderzoek?
% 5. Wat betekenen je resultaten? Wat is de relevantie voor het werkveld?
%
% Daarom bestaat een abstract uit volgende componenten:
%
% - inleiding + kaderen thema
% - probleemstelling
% - (centrale) onderzoeksvraag
% - onderzoeksdoelstelling
% - methodologie
% - resultaten (beperk tot de belangrijkste, relevant voor de onderzoeksvraag)
% - conclusies, aanbevelingen, beperkingen
%
% LET OP! Een samenvatting is GEEN voorwoord!

%%---------- Nederlandse samenvatting -----------------------------------------
%
% TODO: Als je je bachelorproef in het Engels schrijft, moet je eerst een
% Nederlandse samenvatting invoegen. Haal daarvoor onderstaande code uit
% commentaar.
% Wie zijn bachelorproef in het Nederlands schrijft, kan dit negeren, de inhoud
% wordt niet in het document ingevoegd.

\IfLanguageName{english}{%
\selectlanguage{dutch}
\chapter*{Samenvatting}
\lipsum[1-4]
\selectlanguage{english}
}{}

%%---------- Samenvatting -----------------------------------------------------
% De samenvatting in de hoofdtaal van het document

\chapter*{\IfLanguageName{dutch}{Samenvatting}{Abstract}}

Deze bachelorproef behandelt de automatisatie van het onboardingproces binnen de IT-infrastructuur van het gemeentebestuur van Berlare. In veel organisaties verloopt de onboarding van nieuwe medewerkers nog deels manueel, wat leidt tot inefficiëntie, een verhoogde kans op fouten en een langere doorlooptijd voordat medewerkers operationeel inzetbaar zijn.\newline

De probleemstelling van dit onderzoek vertrekt vanuit de nood aan een efficiënter en betrouwbaarder onboardingproces dat vertrekt vanuit gegevens van de personeelsdienst. De centrale onderzoeksvraag luidt: Op welke manier kan een geautomatiseerd systeem ontwikkeld worden dat, op basis van gegevens uit de personeelsdienst, een volledige en correcte onboarding kan uitvoeren binnen de IT-infrastructuur van het gemeentebestuur van Berlare? Het doel van deze bachelorproef is het ontwerpen en implementeren van een dergelijk systeem dat het onboardingproces optimaliseert.\newline

Voor de uitvoering van het onderzoek werd eerst een analyse gemaakt van het bestaande onboardingproces en de IT-omgeving. Vervolgens werd een oplossing ontworpen en technisch uitgewerkt, waarbij gebruik werd gemaakt van automatisatiescripts en integraties met bestaande systemen om gebruikersaccounts en configuraties automatisch aan te maken.\newline

De resultaten tonen aan dat het ontwikkelde systeem in staat is om op een efficiënte en consistente manier nieuwe gebruikers aan te maken en te configureren. Dit leidt tot een aanzienlijke tijdsbesparing en een vermindering van menselijke fouten. Tegelijk werd vastgesteld dat de kwaliteit van de inputgegevens een cruciale factor blijft voor de correcte werking van het systeem.\newline

Uit dit onderzoek kan geconcludeerd worden dat automatisatie een belangrijke meerwaarde biedt voor het onboardingproces binnen lokale besturen. Het systeem verhoogt de efficiëntie en betrouwbaarheid van de IT-dienst en vormt een basis voor verdere uitbreidingen, zoals offboarding of bijkomende integraties. Een beperking van het systeem is de afhankelijkheid van correcte brondata en de nood aan regelmatig onderhoud en opvolging.\newline
